\clearpage\begingroup
\SourceLeaf{325}\HistoricalHead{320}{L}{PRAKTISCHE UND SPHÄRISCHE ASTRONOMIE. BRIEFWECHSEL.}
\begin{center}\textsc{Bessel} an \textsc{Olbers}, 11. Mai 1820\begingroup\renewcommand{\thefootnote}{*)}\footnotemark\renewcommand{\thefootnote}{[*)}\footnotetext{Die Briefe von \textsc{Bessel} an \textsc{Olbers} konnten beim Abdruck nicht mit den Originalen verglichen werden, weil ihr Verbleib, sowie der des ganzen \textsc{Olbers}schen Nachlasses nicht zu ermitteln war.]}\endgroup.\end{center}
\noindent\hspace*{\parindent}\{. . . . . . Ich habe mehrere Prüfungen [mit dem \textsc{Reichenbach}schen Meridiankreise] unter den Händen, unter denen die mit dem Wasserhorizonte die langwierigste werden wird, indem sehr windstilles Wetter dazu erfordert wird, und Sterne, die in dazu schicklichen Höhen culminiren, nicht sehr häufig sind. Inzwischen habe ich etwa 20 Culminationen beobachtet, theils gegen Süden theils gegen Norden, wodurch ich den Einfluss der Biegung vergrössere, indem ich eine Entfernung auf verschiedenen Seiten des Zeniths sowohl direkt als vom Wasserhorizonte beobachte. Sobald die Anzahl der Beobachtungen hinreichend ist, werde ich die Biegung herauszubringen suchen und sie dann mit der auf eine andere Weise gefundenen vergleichen. Ich habe nämlich eine Beobachtungsreihe über den Polarstern angefangen, welche mir die Höhen desselben ganz unabhängig, nicht nur von der Biegung, sondern auch von mehreren anderen, wenigstens möglichen, Fehlern, worüber ich Ihnen ein andermal schreiben werde, gibt: nämlich die Messung seiner doppelten Höhe oder die Entfernung des reflectirten Bildes vom Sterne selbst: dadurch erhalte ich die Polhöhe und diese verglichen mit der durch Umlegung des Instrumentes erhaltenen, gibt die Biegung. . . . . . .\}

\begin{center}\textsc{Gauss} an \textsc{Bessel}, 12. Mai 1820.\end{center}
\noindent\hspace*{\parindent}. . . . . . Ganz unveränderlich scheint der Collimationsfehler[\begingroup\renewcommand{\thefootnote}{**)}\footnotemark\renewcommand{\thefootnote}{[**)}\footnotetext{Indexfehler des Kreises.]}\endgroup\ am \textsc{Reichenbach}schen Meridiankreise] doch nicht zu sein. Ich habe seit dem 21. Febr[uar] bis jetzt erst zweimal umgelegt, in der ersten Periode war der Ort des Pols so gut wie ganz constant, in der dritten aber anfangs merklich kleiner, und jetzt scheint er eher grösser zu sein (in der ersten Periode im Mittel \(321^\circ29'31''\,87\)), wie folgendes Tableau zeigt, wo alle 12 Stunden von einander entfernten Beob[achtungen] eines und desselben Sterns (und zwar, bis auf genauere Reduction, ohne Rücksicht auf die 6 stündige Declinationsänderung) aufgenommen sind. (Sie sehen zugleich daraus, dass ich zuweilen Sterne von der \(5^{\mathrm{ten}}\) Grösse im Mittag beobachte.)
\clearpage\SourceLeaf{326}\HistoricalHead{321}{R}{BEOBACHTUNGEN AM REICHENBACHSCHEN MERIDIANKREISE.}\HistoricalFoot{XI1.}{41}
\begin{center}\begingroup\fontsize{9}{11}\selectfont\setlength{\tabcolsep}{2.3pt}
\begin{tabular}{@{}rlr|rlr@{}}
April 13.&\(\zeta\) Cephei&\(321^\circ29'29\overset{\prime\prime}{,}36\)&April 18.&\(\delta\) Cephei&\(321^\circ29'30\overset{\prime\prime}{,}23\)\\
&\(\gamma\) Cassiopeiae&\(30{,}42\)&&\(\iota\) Cephei&\(29{,}99\)\\
&Nordst[ern]&\(30{,}40\)&&\(\kappa\) Cassiopeiae&\(29{,}83\)\\
&\(\delta\) Cassiop[eiae]&\(29{,}53\)&&Nordstern&\(29{,}49\)\\
&\(\varepsilon\) Cassiop[eiae]&\(29{,}43\)&» 19.&\(\beta\) Cassiop[eiae]&\(31{,}25\)\\
&\(\beta\) Ursae min.&\(30{,}42\)&» 21.&Nordstern&\(30{,}21\)\\
&\(\zeta\) Cephei&\(28{,}74\)&» 22.&\(\gamma\) Cephei&\(30{,}40\)\\
&\(\iota\) Cephei&\(30{,}46\)&» 24.&\(\varepsilon\) Cephei&\(31{,}57\)\\
&\(\alpha\) Ursae mai.&\(29{,}67\)&» 25.&\(\alpha\) Ursae&\(31{,}23\)\\
» 15.&\(\kappa\) Draconis&\(30{,}27\)&» 26.&\(\iota\) Cephei&\(31{,}81\)\\
&\(\omega\) Cephei (5 Gr.)&\(29{,}53\)&&\(\lambda\) Draconis&\(30{,}54\)\\
&Nordstern&\(31{,}17\)::&Mai 10.&\(\lambda\) Draconis\begingroup\renewcommand{\thefootnote}{*)}\raisebox{.65ex}{\footnotemark}\endgroup&\(31{,}70\)\\
» 17.&\(\alpha\) Ursae mai.&\(30{,}53\)&&\(\gamma\) Cephei&\(31{,}79\)\\
&Nordstern&\(29{,}58\)&» 11.&\(\gamma^2\) Ursae min.&\(32{,}54\)\\
» 18.&Nordstern&\(30{,}82\)&&\(\beta\) Ursae min.&\(32{,}84\)
\end{tabular}\par\endgroup\end{center}
\begingroup\renewcommand{\thefootnote}{*)}\footnotetext{Dazwischen in den Vorm[ittags]st[unden] fast immer ungünstig; auch Abhaltung durch meine leidigen Vorlesungen.}\endgroup

\noindent\hspace*{\parindent}Ich möchte wohl wissen, wie Sie, wenn Sie eine ähnliche Erfahrung machen, es mit der Definitiv-Reduction der Beobachtungen halten wollen. Wenn man recht oft umlegt, also sehr kurze Perioden macht, so möchte es am zweckmässigsten sein, für den Ort des Pols das (gehörige) Mittel aus allen Bestimmungen in jeder Periode zum Grunde zu legen; aber das lässt sich nicht immer thun. Wenn in einer längern Periode der Ort des Pols etwas veränderlich erscheint, ohne dass man doch nachweisen kann, \textls[120]{wann} eine Veränderung eingetreten oder inwiefern sie allmählich erzeugt ist, so scheint \textls[120]{einige Willkürlichkeit} bei der Reduction nicht wohl zu vermeiden zu sein. . . . . . .

Um die Resultate unsrer Kreise zu vergleichen, wünschte ich, dass Sie mir einige Beob[achtungen] \textsc{Maskelyne}scher Sterne[\begingroup\renewcommand{\thefootnote}{**)}\footnotemark\renewcommand{\thefootnote}{[**)}\footnotetext{\textsc{Maskelyne} hatte 1805 seinen Katalog von 36 Fundamentalsternen herausgegeben. Die Felder in der folgenden Tabelle enthalten die Anzahl der Beob. dieser 36 Sterne.]}\endgroup] mittheilten. Ich würde gleich Ihnen mehrere schicken, wenn ich wüsste, welche Sie vorzüglich beobachtet haben. Damit Sie Ihrerseits nicht in demselben Fall sind, schreibe ich den Poll[\begingroup\renewcommand{\thefootnote}{***)}\footnotemark\renewcommand{\thefootnote}{[***)}\footnotetext{Englische Bezeichnung für »Register«.]}\endgroup] hier her, wie er heute steht
\clearpage\SourceLeaf{327}\HistoricalHead{322}{L}{PRAKTISCHE UND SPHÄRISCHE ASTRONOMIE. BRIEFWECHSEL.}
\begin{center}\begingroup\setlength{\tabcolsep}{5pt}
\begin{tabular}{r|r|r|r|r|r|r|r|r@{\qquad}lr}
4&12&18&6&4&—&—&—&5&d. i. \(\gamma\) Pegasi&4 mal\\
3&4&11&6&2&1&—&—&—&\(\alpha\) Arietis&3 »\\
4&11&12&4&2&—&—&—&3&\(\alpha\) Ceti&4 »\\
10&9&11&2&1&—&2&2&4&\(\alpha\) Tauri&10 » etc.\\\cline{1-9}
\multicolumn{9}{l}{Summa 153 Beob[achtungen].}&&
\end{tabular}\par\endgroup\end{center}

\noindent\hspace*{\parindent}Auf gut Glück wähle ich den Procyon aus; scheinbare Z[enith]d[istanz] bloss von Refraction nach Ihren Tafeln befreit
\begin{center}\begingroup\setlength{\tabcolsep}{4pt}
\begin{tabular}{rlr}
\multicolumn{3}{c}{Erste Periode, Ort des Pols \(321^\circ29'31\overset{\prime\prime}{,}87\)}\\
52&d. i. Februar 22.&\(45^\circ48'49\overset{\prime\prime}{,}68\)\\
58&&\(\phantom{45^\circ48'}49{,}23\)\\[.7ex]
\multicolumn{3}{c}{Zweite Periode, Ort des Pols \(38^\circ25'53\overset{\prime\prime}{,}54\)}\\
85&&\(314^\circ\ 6'37\overset{\prime\prime}{,}10\)\\
88&&\(\phantom{314^\circ\ 6'}36{,}06\)\\
89&&\(\phantom{314^\circ\ 6'}37{,}29\)\\
90&&\(\phantom{314^\circ\ 6'}36{,}42\)\\[.7ex]
\multicolumn{3}{c}{Dritte Periode, (Data zur Best[immung] des Pols: oben)}\\
106&&\(45^\circ48'46\overset{\prime\prime}{,}61\)\\
107&&\(\phantom{45^\circ48'}46{,}32\)\\
108&&\(\phantom{45^\circ48'}47{,}42\)\\
109&&\(\phantom{45^\circ48'}47{,}93\)\\
117&&\(\phantom{45^\circ48'}47{,}19\)\\
130&&\(\phantom{45^\circ48'}46{,}59.\)
\end{tabular}\par\endgroup\end{center}

\noindent\hspace*{\parindent}Ueber meine Polhöhe wage ich noch nichts zu bestimmen; dürfte man den Ort des Pols der 1. und 2. Periode verbinden, so wäre sie \(51^\circ31'49\overset{\prime\prime}{,}16\); allein ich habe Ursache zu glauben, dass der Collimationsfehler in den ersten Tagen nach der ersten Umlegung sich etwas geändert, das Wetter erlaubte damals nicht gleich, Circumpolarsterne zu beobachten. Beim zweiten Umlegen war es günstiger und ich erhielt von 13 Sternen, die ich vorher beobachtet hatte, die unmittelbar folgende entgegengesetzte Culmination (dies scheint mir im allgemeinen die sicherste Art zu sein) allein leider war damals dazwischen ein kleines Derangement des Netzes (vielleicht durch ein Insect oder sonst etwas darauf gefallenes) gewesen; übrigens gaben diese 13 Sterne
\clearpage\SourceLeaf{328}\HistoricalHead{323}{R}{BEOBACHTUNGEN AM REICHENBACHSCHEN MERIDIANKREISE.}
\noindent im Mittel fast genau \(51^\circ31'48\overset{\prime\prime}{,}00\). Versteht sich, dass dies noch etwanige Flexion \textls[120]{oder alle sonstige} Einwirkung der Schwere involvirt, in Beziehung auf welche \textls[120]{jedes} Instrument als Individuum betrachtet und unerlässlich durch directe Beobachtung der Einfluss bestimmt werden muss (welches übrigens keine Eile hat). Ich habe erst neulich und zwar gerade den Abend vor Empfang Ihres Briefs die erste Beobachtung des Nordsterns auf einem Wasserhorizont gemacht, welche in Zukunft öfters wiederholt werden soll. Jene gab
\begin{center}\begingroup\fontsize{8.5}{11}\selectfont\setlength{\tabcolsep}{2pt}
\begin{tabular}{@{}l@{\quad}l@{}}
Mai 10. oder \(130^{\mathrm{d}}12^{\mathrm{h}}56^{\mathrm{m}}\)&\\
\raisebox{.5\normalbaselineskip}{\(\left.\begin{array}{@{}rrrr@{}}\text{direct}&319^\circ&50'&20\overset{\prime\prime}{,}51\\\text{reflect[irt]}&220&5&3{,}99\end{array}\right\}\)}&\begin{tabular}{@{}rrrrl@{}}W[ahre] Z[enith]—D[istanz]&\(40^\circ\)&\(7'\)&\(21\overset{\prime\prime}{,}74\)&(flüchtig\\Decl[ination]&88&20&\(51{,}05\)&berechnet)\\\cline{1-4}Polhöhe&51&31&\(47{,}21\)&\(-d\delta.\)\end{tabular}
\end{tabular}\par\endgroup\end{center}

\begin{center}\textsc{Bessel} an \textsc{Olbers}, 5. Juni 1820.\end{center}
\noindent\hspace*{\parindent}\{. . . . . . Meine Untersuchungen über die Polhöhe geben jetzt im Mittel unter Voraussetzung der Declination des Polarsterns \(50\overset{\prime\prime}{,}18\); ohne Wasserhorizont \(51\overset{\prime\prime}{,}9\) Polhöhe. . . . . . .\}

\begin{center}\textsc{Bessel} an \textsc{Gauss}, 1. Juni 1820.\end{center}
\noindent\hspace*{\parindent}\{Unter dem herzlichsten Danke für die vielen lehrreichen Mittheilungen Ihres Briefes vom 12. Mai überschicke ich Ihnen hier einige Beobachtungsreihen von \textsc{Maskelyne}schen Sternen, die nun zeigen werden, wie unsere Kreise sich gegen einander verhalten: . . . . . .\}

[folgen Beobachtungen von \(\alpha\) Aurigae, \(\beta\) Tauri, \(\alpha\) Canis mai., \(\alpha\) Geminorum, \(\alpha\) Canis min., \(\beta\) Geminorum, von März 6. bis Mai 25.]
\clearpage\endgroup
